\chapter{The low individual degree test}
\label{chap:test}
\label{sec:the-test}

\section{The game and its strategies}

\begin{definition}[Roles]\label{def:roles}
	\lean{MIPStarRE.LDT.Role}
	\leanok
	The low individual degree test is played between two provers and a
	verifier.  A role is an element of $\{\mathrm A,\mathrm B\}$.  For a role
	$r$, write $\overline r$ for the other role.
\end{definition}

\begin{definition}[Low individual degree test]\label{def:low-individual-degree-test}
	\lean{MIPStarRE.LDT.ProjStrat.lowIndividualDegreeFailureProbability, MIPStarRE.LDT.ProjStrat.PassesLowIndividualDegreeTest}
	\leanok
	\uses{def:roles}
	Let $m\ge1$ and $d\ge0$ be integers and let $q$ be a prime power.  The
	$(m,q,d)$ low individual degree test performs one of the following three
	subtests, each with probability $\frac13$.
	\begin{enumerate}
		\item Axis-parallel lines test: choose a uniformly random role $r$, a
		      uniformly random point $u\in\F_q^m$, and a uniformly random
		      direction $i\in\{1,\ldots,m\}$, and let
		      $\ell=\{u+te_i:t\in\F_q\}$.  Player $r$ receives $\ell$ and
		      returns a univariate degree-$d$ polynomial $f\colon\ell\to\F_q$;
		      Player $\overline r$ receives $u$ and returns $a\in\F_q$.  The
		      verifier accepts if $f(u)=a$.
		\item Self-consistency test: both players receive the same uniformly
		      random point $u$ and must return the same field element.
		\item Diagonal lines test: choose a uniformly random role $r$, a
		      uniformly random point $u$, a uniformly random index
		      $i\in\{1,\ldots,m\}$, and a uniformly random $v\in\F_q^m$ whose last
		      $m-i$ coordinates vanish, and let $\ell=\{u+tv:t\in\F_q\}$.  When
		      $v=0$ this is the singleton $\{u\}$, which is a permitted
		      degenerate line.  Player $r$ receives $\ell$ and returns a
		      degree-$md$ polynomial $f\colon\ell\to\F_q$; Player $\overline r$
		      receives $u$ and returns $a\in\F_q$.  The verifier accepts if
		      $f(u)=a$.
	\end{enumerate}
	A strategy passes the test with probability at least $1-\eps$ if its
	failure probability, the average of the three subtest failure
	probabilities, is at most $\eps$.
\end{definition}

\begin{definition}[Symmetric projective strategy]\label{def:symmetric-projective-strategy}
	\lean{MIPStarRE.LDT.SymStrat, MIPStarRE.LDT.AnswerSymStrat}
	\leanok
	\uses{def:low-individual-degree-test}
	A symmetric projective strategy for the $(m,q,d)$ low individual degree
	test is a tuple $(\psi,A,B,L)$, where $\calH$ is finite-dimensional,
	$\ket\psi\in\calH\ot\calH$ is a permutation-invariant state, and
	\begin{itemize}
		\item $A=\{A^u_a\}$ contains, for each $u\in\F_q^m$, a projective
		      measurement $A^u$ on $\calH$ with outcomes $a\in\F_q$;
		\item $B=\{B^\ell_f\}$ contains, for each axis-parallel line $\ell$, a
		      projective measurement $B^\ell$ with outcomes the univariate
		      degree-$d$ polynomials $f\colon\ell\to\F_q$;
		\item $L=\{L^\ell_f\}$ contains, for each line $\ell$, a projective
		      measurement $L^\ell$ whose outcomes are functions $f\colon\ell\to\F_q$.
	\end{itemize}
	Allowing every function as a diagonal-line answer is a harmless
	relaxation: an ordinary strategy with degree-$md$ answers becomes one with
	function answers by coarse-graining polynomial labels that induce the same
	function on $\ell$, which preserves projectivity and the acceptance
	probability.  Unlike the degree-$md$ answer alphabet, the relaxed alphabet
	is closed under restriction to an affine slice.  The formalization records
	both forms and converts the first into the second.
\end{definition}

\begin{definition}[General projective strategy]\label{def:projective-strategy}
	\lean{MIPStarRE.LDT.ProjStrat}
	\leanok
	\uses{def:low-individual-degree-test}
	A projective strategy is a tuple
	$(\psi,A^{\mathrm A},B^{\mathrm A},L^{\mathrm A},A^{\mathrm B},B^{\mathrm B},L^{\mathrm B})$,
	where $\calH_{\mathrm A}$ and $\calH_{\mathrm B}$ are finite-dimensional,
	$\ket\psi\in\calH_{\mathrm A}\ot\calH_{\mathrm B}$ is a state, and for each
	$w\in\{\mathrm A,\mathrm B\}$ the families $A^w$, $B^w$, $L^w$ are point,
	axis-parallel line, and diagonal-line projective measurements on
	$\calH_w$ as in Definition~\ref{def:symmetric-projective-strategy}.
\end{definition}

\begin{remark}\label{rem:strategy-convention}
	Throughout, a strategy is a projective strategy and a symmetric strategy is
	a symmetric projective strategy.  Most of the argument concerns symmetric
	strategies; the general case is reduced to it in the proof of
	Theorem~\ref{thm:main-formal}.
\end{remark}

\begin{definition}[Good strategy]\label{def:good-strategy}
	\lean{MIPStarRE.LDT.SymStrat.IsGood, MIPStarRE.LDT.AnswerSymStrat.IsGood}
	\leanok
	\uses{def:symmetric-projective-strategy}
	A symmetric strategy is $(\eps,\delta,\gamma)$-good if it passes the
	axis-parallel lines test with probability at least $1-\eps$, the
	self-consistency test with probability at least $1-\delta$, and the
	diagonal lines test with probability at least $1-\gamma$.
\end{definition}

\begin{remark}[Good-strategy characterization]\label{rem:good-strat-characterization}
	In the notation of Chapter~\ref{chap:preliminaries}, a symmetric strategy
	is $(\eps,\delta,\gamma)$-good exactly when, for $\ell$ and $u$ distributed
	as in the axis-parallel lines test,
	\[
		A^u_a\ot I\simeq_\eps I\ot B^\ell_{[f(u)=a]},
		\qquad
		A^u_a\ot I\simeq_\delta I\ot A^u_a,
	\]
	and, for $\ell$ and $u$ distributed as in the diagonal lines test,
	$A^u_a\ot I\simeq_\gamma I\ot L^\ell_{[f(u)=a]}$.
\end{remark}

\begin{definition}[Last-direction lines]\label{not:conditioned-on-last-direction}
	\uses{def:symmetric-projective-strategy}
	Let $(\psi,A,B,L)$ be a symmetric strategy for the $(m+1,q,d)$ test.  For
	$u\in\F_q^m$ write $B^u_f$ for $B^\ell_f$ with
	$\ell=\{(u,x):x\in\F_q\}$, and for a function $f$ on this line write $f(x)$
	for $f(u,x)$.
\end{definition}

\begin{definition}[Restricted diagonal lines test]\label{def:restricted-diagonal-lines-test}
	\lean{MIPStarRE.LDT.RestrictedDiagonalSample}
	\leanok
	\uses{def:low-individual-degree-test}
	For $j\in\{1,\ldots,m\}$, the $j$-restricted diagonal lines test is the
	diagonal lines test conditioned on $i=j$.  In the $m$-restricted test one
	samples $u,v\in\F_q^m$ independently and uniformly and asks the line
	$\{u+tv:t\in\F_q\}$.
\end{definition}

\section{The main theorem}

\begin{definition}[The error parameter]\label{def:main-formal-error}
	\lean{MIPStarRE.LDT.Test.simplifiedFinalScale, MIPStarRE.LDT.Test.simplifiedMainFormalError}
	\leanok
	For integers $m\ge1$ and $d\ge0$, a prime power $q$, and $\eps\ge0$, set
	\[
		K_{m,d}=\begin{cases}m^4,&d=0,\\d^2m^4,&d\ge1,\end{cases}
		\qquad
		\nu=21000\,K_{m,d}\bigl(\eps^{1/64}+(d/q)^{1/64}\bigr).
	\]
\end{definition}

\begin{theorem}[Quantum soundness of the low individual degree test]\label{thm:main-formal}
	\lean{MIPStarRE.LDT.Test.mainFormal}
	\leanok
	\uses{def:projective-strategy, def:low-individual-degree-test, def:polymeasurements, def:simeq, def:main-formal-error}
	Let $m\ge1$ and let
	$(\psi,A^{\mathrm A},B^{\mathrm A},L^{\mathrm A},A^{\mathrm B},B^{\mathrm B},L^{\mathrm B})$
	be a projective strategy which passes the $(m,q,d)$ low individual degree
	test with probability at least $1-\eps$.  Let $\nu$ be as in
	Definition~\ref{def:main-formal-error}.  Then there exist projective
	measurements $G^{\mathrm A},G^{\mathrm B}\in\polymeas mqd$ such that
	\begin{enumerate}
		\item (consistency with $A$) on average over $u\sim\F_q^m$,
		      \[
			      A^{\mathrm A,u}_a\ot I\simeq_\nu I\ot G^{\mathrm B}_{[g(u)=a]},
			      \qquad
			      G^{\mathrm A}_{[g(u)=a]}\ot I\simeq_\nu I\ot A^{\mathrm B,u}_a;
		      \]
		\item (self-consistency)
		      $G^{\mathrm A}_g\ot I\simeq_\nu I\ot G^{\mathrm B}_g$.
	\end{enumerate}
\end{theorem}

\begin{proof}
	\leanok
	\uses{lem:main-formal-saturated, lem:main-formal-small-error}
	If $\nu\ge1$, apply Lemma~\ref{lem:main-formal-saturated} below.
	Otherwise apply Lemma~\ref{lem:main-formal-small-error}, which is proved
	in Section~\ref{sec:proof-of-main-formal} after the main induction.
\end{proof}



\begin{lemma}[Saturated error]\label{lem:main-formal-saturated}
	\lean{MIPStarRE.LDT.Test.mainFormalConclusion_of_one_le}
	\leanok
	\uses{def:projective-strategy, def:polymeasurements, def:simeq}
	This lemma is a formalization-only auxiliary result supporting the first
	paragraph of the proof of Theorem~\ref{thm:main-formal}.  For every
	projective strategy and every $\nu\ge1$, the conclusions of
	Theorem~\ref{thm:main-formal} hold with the projective measurements
	supported on a fixed polynomial $g_0\in\polyfunc mqd$.
\end{lemma}
\begin{proof}
	\leanok
	The self-inconsistency of these measurements vanishes, and every point
	inconsistency on a normalized state is at most one.
\end{proof}

The original paper states the theorem with an auxiliary sampling parameter
$k$.  That statement follows from Theorem~\ref{thm:main-formal}.

\begin{definition}[The sampling-parameter error]\label{def:main-formal-sampling-error}
	\lean{MIPStarRE.LDT.Test.mainFormalError}
	\leanok
	For an integer $k$, set
	\[
		\nu_k=100000\,k^2m^4\Bigl(\eps^{1/40000}+(d/q)^{1/40000}
		+e^{-k/(2560000m^2)}\Bigr).
	\]
\end{definition}

\begin{lemma}[Comparison of the error parameters]\label{lem:main-formal-error-comparison}
	\lean{MIPStarRE.LDT.Test.simplifiedMainFormalError_le_mainFormalError}
	\leanok
	\uses{def:main-formal-error, def:main-formal-sampling-error}
	Let $\eps\ge0$ and let $k\ge md$ be a positive integer.  If $\nu_k<1$, then
	$\nu\le\nu_k$.
\end{lemma}
\begin{proof}
	\leanok
	Since $k,m\ge1$, the prefactor $100000k^2m^4$ is at least one, so
	$\nu_k<1$ forces $\eps^{1/40000}<1$ and $(d/q)^{1/40000}<1$, hence
	$\eps,d/q\in[0,1]$.  On $[0,1]$ the power $x^{1/64}$ is at most
	$x^{1/40000}$.  Finally $K_{m,d}\le k^2m^4$: for $d=0$ because $k\ge1$, and
	for $d\ge1$ because $d\le md\le k$.  Hence
	$21000K_{m,d}(\eps^{1/64}+(d/q)^{1/64})$ is at most
	$100000k^2m^4(\eps^{1/40000}+(d/q)^{1/40000})\le\nu_k$.
\end{proof}

\begin{corollary}[The sampling-parameter statement]\label{cor:main-formal-sampling}
	\lean{MIPStarRE.LDT.Test.mainFormalWithK}
	\leanok
	\uses{def:projective-strategy, def:low-individual-degree-test, def:polymeasurements, def:simeq, def:main-formal-sampling-error}
	Let the strategy be as in Theorem~\ref{thm:main-formal}, and let $k$ be a
	positive integer with $k\ge md$.  Then the conclusions of
	Theorem~\ref{thm:main-formal} hold with $\nu$ replaced by $\nu_k$.  This
	is the statement of the main theorem of
	\cite{Ji2020LowIndividualDegree}, with the additional condition $k>0$: at
	$k=0$, which the printed hypothesis permits when $d=0$, the printed error
	vanishes \cite{gap:issue-422-main-formal-zero-k-boundary}.
\end{corollary}
\begin{proof}
	\leanok
	\uses{thm:main-formal, lem:main-formal-saturated, lem:main-formal-error-comparison}
	If $\nu_k\ge1$, apply Lemma~\ref{lem:main-formal-saturated}.  Otherwise
	Lemma~\ref{lem:main-formal-error-comparison} gives $\nu\le\nu_k$, and the
	conclusions of Theorem~\ref{thm:main-formal} weaken from $\nu$ to $\nu_k$.
\end{proof}
